% Options for packages loaded elsewhere
\PassOptionsToPackage{unicode}{hyperref}
\PassOptionsToPackage{hyphens}{url}
\documentclass[
]{article}
\usepackage{xcolor}
\usepackage{amsmath,amssymb}
\setcounter{secnumdepth}{-\maxdimen} % remove section numbering
\usepackage{iftex}
\ifPDFTeX
  \usepackage[T1]{fontenc}
  \usepackage[utf8]{inputenc}
  \usepackage{textcomp} % provide euro and other symbols
\else % if luatex or xetex
  \usepackage{unicode-math} % this also loads fontspec
  \defaultfontfeatures{Scale=MatchLowercase}
  \defaultfontfeatures[\rmfamily]{Ligatures=TeX,Scale=1}
\fi
\usepackage{lmodern}
\ifPDFTeX\else
  % xetex/luatex font selection
\fi
% Use upquote if available, for straight quotes in verbatim environments
\IfFileExists{upquote.sty}{\usepackage{upquote}}{}
\IfFileExists{microtype.sty}{% use microtype if available
  \usepackage[]{microtype}
  \UseMicrotypeSet[protrusion]{basicmath} % disable protrusion for tt fonts
}{}
\makeatletter
\@ifundefined{KOMAClassName}{% if non-KOMA class
  \IfFileExists{parskip.sty}{%
    \usepackage{parskip}
  }{% else
    \setlength{\parindent}{0pt}
    \setlength{\parskip}{6pt plus 2pt minus 1pt}}
}{% if KOMA class
  \KOMAoptions{parskip=half}}
\makeatother
\usepackage{longtable,booktabs,array}
\newcounter{none} % for unnumbered tables
\usepackage{calc} % for calculating minipage widths
% Correct order of tables after \paragraph or \subparagraph
\usepackage{etoolbox}
\makeatletter
\patchcmd\longtable{\par}{\if@noskipsec\mbox{}\fi\par}{}{}
\makeatother
% Allow footnotes in longtable head/foot
\IfFileExists{footnotehyper.sty}{\usepackage{footnotehyper}}{\usepackage{footnote}}
\makesavenoteenv{longtable}
\setlength{\emergencystretch}{3em} % prevent overfull lines
\providecommand{\tightlist}{%
  \setlength{\itemsep}{0pt}\setlength{\parskip}{0pt}}
\usepackage{bookmark}
\IfFileExists{xurl.sty}{\usepackage{xurl}}{} % add URL line breaks if available
\urlstyle{same}
\hypersetup{
  hidelinks,
  pdfcreator={LaTeX via pandoc}}

\author{}
\date{}

\begin{document}

\section{Probability Warm-Up: Intersection, Union, and
Conditional}\label{probability-warm-up-intersection-union-and-conditional}

\emph{Companion handout --- complete before the Bayes' Theorem lecture}

Name: \_\_\_\_\_\_\_\_\_\_\_\_\_\_\_\_\_\_\_\_\_\_ Date:
\_\_\_\_\_\_\_\_\_\_\_\_\_\_

\begin{center}\rule{0.5\linewidth}{0.5pt}\end{center}

\subsection{0. The vocabulary}\label{the-vocabulary}

An \textbf{experiment} is anything with an uncertain outcome: rolling a
die, testing a patient, drawing a card.

The \textbf{sample space} \(S\) is the set of all possible outcomes. For
one die roll, \(S = \{1,2,3,4,5,6\}\).

An \textbf{event} is a subset of the sample space --- something that
either happens or doesn't. ``The roll is even'' is the event
\(A = \{2,4,6\}\).

When every outcome is equally likely:

\[P(A) = \frac{\text{number of outcomes in } A}{\text{number of outcomes in } S}\]

So \(P(\text{even}) = 3/6 = 0.5\).

Two facts we will lean on constantly:

\begin{itemize}
\tightlist
\item
  \(0 \le P(A) \le 1\) always.
\item
  \(P(\bar{A}) = 1 - P(A)\), where \(\bar{A}\) (``not \(A\)'') is the
  \textbf{complement} of \(A\).
\end{itemize}

\begin{quote}
\textbf{Quick check.} A single die is rolled. Write out the event \(B\)
= ``the roll is greater than 4'' and compute \(P(B)\) and
\(P(\bar{B})\).

\(B = \{\rule{2.5cm}{0.15pt}\}\) \(P(B) = \rule{1.5cm}{0.15pt}\)
\(P(\bar{B}) = \rule{1.5cm}{0.15pt}\)
\end{quote}

\begin{center}\rule{0.5\linewidth}{0.5pt}\end{center}

\subsection{\texorpdfstring{1. Intersection ---
``\textbf{and}''}{1. Intersection --- ``and''}}\label{intersection-and}

\(A \cap B\) is the event that \textbf{both} \(A\) and \(B\) happen.
Read the symbol as ``and.''

\textbf{Example.} Draw one card from a standard 52-card deck. Let

\begin{itemize}
\tightlist
\item
  \(A\) = the card is a heart (\(13\) cards)
\item
  \(B\) = the card is a face card --- J, Q, K (\(12\) cards)
\end{itemize}

\(A \cap B\) = the card is a \emph{heart face card}. There are exactly 3
of those (J♥, Q♥, K♥), so

\[P(A \cap B) = \frac{3}{52} \approx 0.058\]

Notice that \(P(A \cap B)\) is smaller than both \(P(A) = 13/52\) and
\(P(B) = 12/52\). Adding a requirement can only shrink an event --- it
never grows it. This is worth saying out loud, because it is the sanity
check that catches most arithmetic errors later.

If \(A \cap B\) is empty --- the two events cannot co-occur --- we call
them \textbf{mutually exclusive} and \(P(A \cap B) = 0\). ``The card is
a heart'' and ``the card is a spade'' are mutually exclusive.

\begin{quote}
\textbf{Exercise 1.} One card is drawn. Let \(A\) = ``the card is red''
and \(B\) = ``the card is an ace.''

\begin{enumerate}
\def\labelenumi{(\alph{enumi})}
\tightlist
\item
  How many cards are in \(A \cap B\)? \_\_\_\_\_\_\_
\item
  \(P(A \cap B)=\) \_\_\_\_\_\_\_
\item
  Name two events in this deck that are mutually exclusive.
\end{enumerate}
\end{quote}

\begin{center}\rule{0.5\linewidth}{0.5pt}\end{center}

\subsection{\texorpdfstring{2. Union ---
``\textbf{or}''}{2. Union --- ``or''}}\label{union-or}

\(A \cup B\) is the event that \(A\) happens, or \(B\) happens, or both.
Read the symbol as ``or.''

The tempting move is to write \(P(A \cup B) = P(A) + P(B)\). That is
wrong whenever the events overlap, because everything in the overlap
gets counted twice --- once in \(P(A)\) and again in \(P(B)\). The fix
is to subtract the double-count:

\[P(A \cup B) = P(A) + P(B) - P(A \cap B)\]

\textbf{Example.} Same deck, \(A\) = heart, \(B\) = face card.

\[P(A \cup B) = \frac{13}{52} + \frac{12}{52} - \frac{3}{52} = \frac{22}{52} \approx 0.423\]

Sanity check by counting directly: 13 hearts, plus the 9 face cards that
aren't hearts, is 22 cards. ✓

When the events \emph{are} mutually exclusive, \(P(A \cap B) = 0\) and
the formula collapses back to plain addition. That's the special case,
not the general rule.

\begin{quote}
\textbf{Exercise 2.} In a class of 40 students, 22 have written Python
before, 15 have taken a statistics course, and 9 have done both.

\begin{enumerate}
\def\labelenumi{(\alph{enumi})}
\tightlist
\item
  How many have done Python \emph{or} statistics (or both)?
  \_\_\_\_\_\_\_
\item
  \(P(\text{Python} \cup \text{Stats})=\) \_\_\_\_\_\_\_
\item
  How many have done \emph{neither}? \_\_\_\_\_\_\_
\item
  Sketch a Venn diagram and label all four regions with counts.
\end{enumerate}
\end{quote}

\begin{center}\rule{0.5\linewidth}{0.5pt}\end{center}

\subsection{\texorpdfstring{3. Conditional probability ---
``\textbf{given}''}{3. Conditional probability --- ``given''}}\label{conditional-probability-given}

\(P(A \mid B)\) is the probability that \(A\) happens \textbf{given that
we already know \(B\) happened}. Read the bar as ``given.''

Conditioning on \(B\) means throwing away every outcome where \(B\)
didn't happen. \(B\) becomes the new sample space. So we count the
outcomes where both happened, and divide by the outcomes where \(B\)
happened:

\[P(A \mid B) = \frac{P(A \cap B)}{P(B)}\]

\textbf{Example.} Draw a card. What is
\(P(\text{face card} \mid \text{heart})\)?

\[P(B \mid A) = \frac{P(A \cap B)}{P(A)} = \frac{3/52}{13/52} = \frac{3}{13} \approx 0.231\]

Or just reason it out: we know it's a heart, so there are 13
possibilities, and 3 of them are face cards. Same answer, and the
formula is only bookkeeping for that intuition.

\subsubsection{Two things to notice}\label{two-things-to-notice}

\textbf{Order matters.} \(P(A \mid B)\) and \(P(B \mid A)\) are
different questions with different denominators. Above,
\(P(\text{face} \mid \text{heart}) = 3/13\), but
\(P(\text{heart} \mid \text{face}) = \frac{3/52}{12/52} = 3/12 = 1/4\).
Same numerator, different denominator, different answer.
\textbf{Confusing these two is the single most common mistake in all of
probability}, and the entire point of Bayes' Theorem is to convert one
into the other correctly.

\textbf{Sometimes knowing \(B\) tells you nothing.} If
\(P(A \mid B) = P(A)\), the events are \textbf{independent} and the
definition rearranges into \(P(A \cap B) = P(A)P(B)\). Two coin flips
are independent. A card being a heart and being a face card are also
independent --- check it: \(3/13 = P(\text{face} \mid \text{heart})\)
and \(12/52 = 3/13 = P(\text{face})\). Equal, so independent. But an
infection and a test result had better \emph{not} be independent, or the
test is useless.

\begin{quote}
\textbf{Exercise 3.} Roll two fair dice. Let \(A\) = ``the sum is 8''
and \(B\) = ``the first die shows 5.''

\begin{enumerate}
\def\labelenumi{(\alph{enumi})}
\tightlist
\item
  \(P(A)=\) \_\_\_\_\_\_\_ (there are 36 equally likely outcomes)
\item
  \(P(A \cap B)=\) \_\_\_\_\_\_\_
\item
  \(P(A \mid B)=\) \_\_\_\_\_\_\_
\item
  Is \(P(A \mid B)\) bigger or smaller than \(P(A)\)? What does that
  tell you?
\item
  Compute \(P(B \mid A)\) and confirm it differs from \(P(A \mid B)\).
\end{enumerate}
\end{quote}

\begin{center}\rule{0.5\linewidth}{0.5pt}\end{center}

\subsection{4. Putting it together: the two-way
table}\label{putting-it-together-the-two-way-table}

Most real problems arrive as counts in a table, and every quantity above
can be read straight off it.

Out of \textbf{1,000 people}, 1\% carry an infection. A test is applied
to all of them; among the infected it comes back positive 97\% of the
time, and among the healthy it also comes back positive 3\% of the time.

{\def\LTcaptype{none} % do not increment counter
\begin{longtable}[]{@{}lrrr@{}}
\toprule\noalign{}
& Test \textbf{+} & Test \textbf{−} & \textbf{Total} \\
\midrule\noalign{}
\endhead
\bottomrule\noalign{}
\endlastfoot
\textbf{Infected} & 9.7 & 0.3 & 10 \\
\textbf{Not infected} & 29.7 & 960.3 & 990 \\
\textbf{Total} & 39.4 & 960.6 & 1,000 \\
\end{longtable}
}

(Decimal people are fine --- these are expected counts, not a
headcount.)

Reading the table:

\begin{itemize}
\tightlist
\item
  \textbf{Intersection:} \(P(I \cap +) = 9.7/1000 = 0.0097\) --- one
  cell, over the grand total.
\item
  \textbf{Marginal:} \(P(+) = 39.4/1000 = 0.0394\) --- a row or column
  total, over the grand total.
\item
  \textbf{Conditional:} \(P(+ \mid I) = 9.7/10 = 0.97\) --- one cell,
  over its \emph{row} total. Conditioning on \(I\) means the
  ``Infected'' row is now the whole world.
\item
  \textbf{The other conditional:}
  \(P(I \mid +) = 9.7/39.4 \approx 0.246\) --- the same cell, over its
  \emph{column} total.
\end{itemize}

\begin{quote}
\textbf{Exercise 4.} Using the table above:

\begin{enumerate}
\def\labelenumi{(\alph{enumi})}
\tightlist
\item
  \(P(\bar{I} \cap +)=\) \_\_\_\_\_\_\_
\item
  \(P(- \mid \bar{I})=\) \_\_\_\_\_\_\_
\item
  \(P(I \cup +)=\) \_\_\_\_\_\_\_ \emph{(use the union formula, then
  verify by counting cells)}
\item
  Are \(I\) and \(+\) independent? Justify with numbers.
\end{enumerate}
\end{quote}

\begin{quote}
\textbf{Exercise 5 --- the punchline.} Compare your answers for
\(P(+ \mid I) = 0.97\) and \(P(I \mid +) \approx 0.246\).

Both describe ``the test and the infection agreeing.'' In one sentence,
explain to a classmate why they are so far apart. Which of the two would
a patient actually want to know?
\end{quote}

\begin{center}\rule{0.5\linewidth}{0.5pt}\end{center}

\subsection{Answer Key}\label{answer-key}

\textbf{Quick check.} \(B = \{5,6\}\), \(P(B) = 2/6 = 1/3\),
\(P(\bar{B}) = 2/3\).

\textbf{Exercise 1.} (a) 2 (A♥, A♦). (b) \(2/52 = 1/26 \approx 0.038\).
(c) e.g.~``heart'' and ``spade,'' or ``ace'' and ``king'' --- any two
that can't hold at once.

\textbf{Exercise 2.} (a) \(22 + 15 - 9 = 28\). (b) \(28/40 = 0.70\). (c)
\(40 - 28 = 12\). (d) Python only 13, both 9, Stats only 6, neither 12
--- the four must sum to 40.

\textbf{Exercise 3.} (a) \(5/36\) --- the sum-8 outcomes are
(2,6),(3,5),(4,4),(5,3),(6,2). (b) \(1/36\) --- only (5,3). (c)
\(\frac{1/36}{6/36} = 1/6\). (d) \(1/6 > 5/36\), so learning the first
die is a 5 makes a sum of 8 \emph{more} likely; the events are not
independent. (e) \(P(B \mid A) = \frac{1/36}{5/36} = 1/5\), which is not
\(1/6\).

\textbf{Exercise 4.} (a) \(29.7/1000 = 0.0297\). (b)
\(960.3/990 \approx 0.970\). (c)
\(P(I) + P(+) - P(I \cap +) = 0.01 + 0.0394 - 0.0097 = 0.0397\); by
cells, \((9.7 + 0.3 + 29.7)/1000 = 0.0397\). ✓ (d)
No.~\(P(+ \mid I) = 0.97\) but \(P(+) = 0.0394\) --- wildly different,
so the test result depends heavily on infection status. (Good! An
independent test would be worthless.)

\textbf{Exercise 5.} The infection is rare, so the 990 healthy people
vastly outnumber the 10 infected ones --- and 3\% of a large group (29.7
false positives) beats 97\% of a tiny group (9.7 true positives). A
positive result is therefore most often a false alarm. The patient wants
\(P(I \mid +)\), but the test's advertised ``97\% accuracy'' reports
\(P(+ \mid I)\). \textbf{Bayes' Theorem is the tool that converts the
number we are given into the number we want.}

\end{document}
